%%%%%%%%%%%%%%%%%%%%%%%%%%%%%%%%%%%%%%%%%%%%%%%%%%%%%%%%%%%%%%%%%%%%%%%%%%
%%%%%%%%%%%%%%%%%%%%%%%%%%%%%%%%%%%%%%%%%%%%%%%%%%%%%%%%%%%%%%%%%%%%%%%%%%
\chapter{Deuxième chapitre}
\label{chap:two}

\epigraph{\itshape Si resté est ce texte, 0/20 tu auras bien mérité}{-- Gava Yoda}

\stopcontents[chapters]
\startcontents[chapters]

\printMyMinitoc

Texte après le mini-plan que vous pouvez mettre ici. Vous pouvez parler du contenu du chapitre et rappelez où nous en sommes dans l'avancement du mémoire.\newline
{\noindent \color{colorEPISEN-MEMOIRE}\textbf{Objectifs}}:
\begin{itemizeEPISEN}
	\item Réviser les bases ...
	\item Analyser ... 
	\item Rappeler ...
\end{itemizeEPISEN}

\fancyhead[LE]{\textsc{\chaptername~\thechapter}}

%%%%%%%%%%%%%%%%%%%%%%%%%%%%%%%%%%%%%%%%%%%%%%%%%%%%%%%%%%%%%%%%%%%%%%%%%%
%%%%%%%%%%%%%%%%%%%%%%%%%%%%%%%%%%%%%%%%%%%%%%%%%%%%%%%%%%%%%%%%%%%%%%%%%%


\section{Titre section 1 (second chapitre)}
\label{sec:two:one}

\lipsum[4] 

\lipsum[5] 

\section{Des codes}
\label{sec:two:two}

\lipsum[6]

\subsection{Java}

\lipsum[7]

\begin{java}
public class {
	...
	public static void main(String arg){
		System.out.println("sdmklfjsdfj dfsdf ");
		/* A loop */
		for(int i=0;i<n;i++)
			if (1=1)
				Toto.add(i)
	}
}
\end{java}

\lipsum[8]

\subsection{Code Python}

\lipsum[9]

On peut aussi mettre une barre devant les codes (exemple ci-dessous) et même
mettre dans des figures (voir Figure~\ref{fig:code1}). Cela fait beaucoup de couleur je trouve mais si vous aimez, ca fonctionne.

\begin{code}
\begin{python}
# Les listes :
maListeDeMeubles = ['table','chaise','frigo']
maListeDeMeubles.sort()  #Tri de la liste

# A loop et avec des maths $x=\int_{i=1}^n \sin{x}$
for unMeuble in maListeDeMeubles:
   print('longueur de la chaîne ', unMeuble, '=', len(unMeuble))

#les listes imbriquées:
laListeDesNombresPairs = [unNombre for unNombre in range(1000) if unNombre % 2 == 0]
print (laListeDesNombresPairs)

#Les dictionnaires :
unAnnuaire = {'Laurent': 6389565, 'Paul': 6356785}
for unNom, x in unAnnuaire.items():
   print("le nom %s a pour numéro de téléphone %d" %(unNom, x))


#Les tuples (n-uplet) : séquence constante
Couleur = ('Rouge', 'Bleu', 'Vert')
print(Couleur[0], Couleur[1],  Couleur[2])

PointDeReference = (100, 200)
print(" x0 = %d   y0 = %d "  %(PointDeReference[0],PointDeReference[1]))
\end{python}
\end{code}

\lipsum[10]

\subsection{Code R (maths)}

\lipsum[11]

\begin{R}
data <- read.table("data.dat", header = FALSE, sep = "")
M <- data.frame(n1, n2, m)
x <- 1:30

fib <- function(n) {
|	if (n < 2)
|	| n
|	else
|	| fib(n - 1) + fib(n - 2)
}

fib(10) # => 55

myfun <- function(S, F) {
|	data <- read.table(F)
|	plot(data$V1, data$V2, type="l")
|	title(S)
}

for(i in 1:length(species)) {
|	data <- read.table(file[i]) # lit les données
|	plot(data$V1, data$V2, type="l")
|	title(species[i]) # ajoute le titre
}

#Calculate CE for each counterparty
Value.A <- data.frame() #MTM value of each contract within cp A
# ...
for(i in 1:length(foo)){
|	if(isTrue(as.character(portfolio_data[i,1])=="A")==TRUE){
|	# ...
}
\end{R}

\lipsum[12]

\subsection{Code source sans formatage}

\lipsum[13]

Exemple d'un Makefile (mais mettez ce que vous voulez, du HTML, des scripts, la sortie
standard de votre shell, {\it etc.} ou du textt+maths)
\begin{codeExample}
\begin{small}
\begin{verbatim}
fast:
	pdflatex memoire.tex

memoire: all

manuscrit: memoire

all:
	pdflatex memoire.tex
	pdflatex memoire.tex

clean:
	rm -f memoire.pdf memoire.aux  memoire-blx.bib memoire.lot memoire.mtc0 	

allClean: clean
	rm -f *.pdf
\end{verbatim}
\end{small}
\end{codeExample}

\lipsum[14]

\lipsum[15]

\begin{figure}[!t]
\begin{code}
\begin{java}
public class {
|	...
|	public static void main(String arg){
|	|	System.out.println("sdmklfjsdfj dfsdf ");
|	|	/* A loop et avec des maths $x=\int_{i=1}^n \sin{x}$ */
|	|	for(int i=0;i<n;i++)
|	|	|	if (1=1)
|	|	|	|	Toto.add(i)
|	}
}
\end{java}
\end{code}
\caption{Un exemple de code Java dans une figure.}
\label{fig:code1}
\end{figure}

Et un exemple d'Algorithme~\ref{algo:level}. Vous pouvez faire une table des algorithme mais je laisse cela en exercice à ceux qui en auraient besoin...

\begin{algorithm}[!t]
\SetAlgoLined
\SetKwInOut{Input}{\color{colorEPISEN-MEMOIRE}input}
\SetKwInOut{Output}{\color{colorEPISEN-MEMOIRE}output}
\SetKwData{indegreeMap}{indegreeMap}
%\SetKwFunction{FindCompress}{FindCompress}
\Input{A boolean circuit $G$}
\Output{The level of each gate $v \in G$}
\BlankLine
inLevelMap $\leftarrow \{ (v\!\Rightarrow\!0) \ \ \forall v\!\in V_{in} \}$ \;
slice $\leftarrow [ v                 \ \ \forall v\!\in V_{in} ]$ \;
accu $\leftarrow \emptyset$ \;
numOfSlices $\leftarrow 0$ \;
\While{slice $\neq \emptyset$}{
	accu $\leftarrow$ accu $\oplus$ slice \;
	newSlice $\leftarrow$ $[]$ \;
	\For{v $\in$ slice}{
		\For{(v',e) $\in \mathbf{edges}(G,v)$}{
			inLevelMap[e] $\leftarrow$ inLevelMap[e]-1 \;
			\If{inLevelMap[e] $\neq 0$}{
				newSlice $\leftarrow$ newSlice $\oplus$ e
			}
		}
	}
	slice $\leftarrow$ newSlice \;
	numOfSlices++
}
\Return accu, numOfSlices-1

%$\{v: d \ for v, d in G.in_degree().items() if d > 0 \}$ \;
%\While{not at end of this document}{
%read current\;
%\eIf{understand}{
%go to next section\;
%current section becomes this one\;
%}{
%go back to the beginning of current section\;
%}
%}
\caption{Building the level of each gate.}
\label{algo:level}
\end{algorithm}

\lipsum[16]